\scanpage{22}
% archive sheet number: 21
\mspage{- 9 -}\hfill n$^{\mathrm{o}}$ 276

\uline{Remarque} - La réciproque du corollaire précédent est fausse ; cf. Exerc.N.

4. - \uline{Modules filtrés complets}.

Soit $(E,E_n)$ un module filtré sur un anneau filtré $(A,A_n)$. Puisque la structure uniforme du groupe topologique $E$ est définie par un seul écart, il est clair que $E$ \uline{est complet si et seulement si toute suite de Cauchy, dans} $E$ \uline{est convergente}. Par ailleurs en vertu de l'inégalité
\[
d(x_p,x_q)\handwritten{\leqslant}\underset{p\handwritten{\leqslant} n\handwritten{\leqslant} q-1}{\mathrm{Max}}\;d(x_n,x_{n+1})\quad,
\]
on voit qu'une suite $(x_n)$ de points de $E$ vérifie le critère de Cauchy si et seulement si la suite $d(x_n,x_{n+1})$ tend vers 0, i.e. si et seulement si la suite de terme général $y_n=x_{n+1}-x_n$ tend vers 0 dans $E$. On en déduit :

PROPOSITION 4 - \uline{Pour que le module filtré} $E$ \uline{soit complet}, \uline{il faut et il suffit que toute série d'éléments de} $E$ \uline{dont le terme général tend vers} 0 \uline{soit convergente}.

\uline{Corollaire} 1- \uline{Soit} $(E,E_n)$ \uline{un module filtré complet}. \uline{Tout sous-module fermé de} $E$, \uline{muni de la filtration induite}, \uline{est complet}. \uline{Tout module quotient de} $E$, \uline{muni de la filtration quotient est complet}.
% slanted
\margin{\handwritten{\unsure{faits généraux|fait général}\par \unsure{pour les top|pour les esp.}\par métrisables}}

\margin{\handwritten{!}}L'assertion concernant les sous-modules \uline{résulte} trivialement de la Proposition. D'autre part, si l'on considère, dans un module quotient $E/F$ une série dont le terme général tend vers 0, il est clair que celle-ci est l'image d'une série dans $E$ dont le terme général tend vers 0, donc d'une série convergente dans $E$ ; l'application canonique $E\handwritten{\longrightarrow} E\handwritten{/}F$ étant continue, $E/F$ est complet d'après la proposition 1.

\uline{Corollaire} 2 - \uline{Soient} $(A,A_n)$ \uline{un anneau filtré et complet}, $E$ \uline{un} $A$-\uline{module de type fini}, \uline{filtré par les sous-groupes} $A_n E$. \uline{Alors} $E$ \uline{est complet}

Ecrivons $E=L\handwritten{/}R$ où $L$ est un $A$-module libre de type fini ; la filtration de $E$ par les $A_n E$ se déduit alors de la filtration analogue de $L$ ;

